% =====================================================================
%  Kevin Setter, "Topological Quantum Field Theory and the Geometric
%  Langlands Correspondence"（Caltech 博士论文, 2013）
%  —— 中文译注本导言区（Lead 独占维护，译者只读）
%
%  编译：xelatex -interaction=nonstopmode -file-line-error main.tex  (×3)
%  本机：MiKTeX，PATH 直接可用；TMP 已设为可写目录。
%  约定：
%   * 编号公式一律 \tag{原文编号}，与原文一一对应；
%   * 图 1.1–6.1 为原文矢量图 400 dpi 截取，图注为译者所译；
%   * % [Tutor Note: ...] 为译者批注，不进入排版，汇总见 译注汇总.md；
%   * 原文 "Chapter N" → 第 N 章；"Appendix A" → 附录 A。
% =====================================================================
\documentclass[11pt,a4paper,openany,UTF8,fontset=none]{ctexbook}

% ---------- 字体（Windows 绝对路径，不依赖 fontconfig） ----------
\setCJKmainfont[Path=C:/Windows/Fonts/,BoldFont=simhei.ttf,ItalicFont=simkai.ttf]{simsun.ttc}
\setCJKsansfont[Path=C:/Windows/Fonts/]{simhei.ttf}
\setCJKmonofont[Path=C:/Windows/Fonts/]{simfang.ttf}

\usepackage{amsmath,amssymb,amsthm,mathtools}
\usepackage{bm}
\usepackage{graphicx}
\usepackage[a4paper,left=2.6cm,right=2.6cm,top=2.7cm,bottom=2.7cm,headsep=0.5cm]{geometry}
\usepackage[dvipsnames]{xcolor}
\usepackage{enumitem}
\usepackage{booktabs}
\usepackage{array}
\usepackage{multirow}
\usepackage{longtable}
\usepackage{multicol}
\usepackage{fancyhdr}
\usepackage{caption}
\usepackage{tikz}
\usetikzlibrary{arrows.meta,positioning,decorations.pathmorphing,calc}
\usepackage[hidelinks,unicode]{hyperref}
\hypersetup{
  pdftitle={拓扑量子场论与几何朗兰兹对应（中文译注本）},
  pdfauthor={Kevin Setter; 中译与批注：本译本},
  pdfsubject={Caltech PhD thesis, 2013}
}

% ---------- 版式 ----------
\linespread{1.18}
\setlength{\emergencystretch}{3em}
\allowdisplaybreaks[1]
\graphicspath{{figures/}{../figures/}}
\captionsetup{font=small,labelsep=quad}
\setcounter{secnumdepth}{3}
\setcounter{tocdepth}{3}

% ---------- 章节标题样式 ----------
\ctexset{
  chapter/format = \Large\bfseries\centering,
  chapter/name   = {第,章},
  section/format = \large\bfseries,
}
\renewcommand{\chaptername}{第}
\renewcommand{\thechapter}{\arabic{chapter}}

% ---------- 页眉页脚 ----------
\pagestyle{fancy}
\fancyhf{}
\fancyhead[L]{\small\itshape Topological Quantum Field Theory and the Geometric Langlands Correspondence}
\fancyhead[R]{\small\thepage}
\renewcommand{\headrulewidth}{0.4pt}
\fancypagestyle{plain}{\fancyhf{}\fancyfoot[C]{\small\thepage}\renewcommand{\headrulewidth}{0pt}}

% ---------- 定理类环境（编号一律由原文手工给出） ----------
% 用法：\begin{theorem}[3.3] …… \end{theorem}  → 排版为「定理 3.3. ……」
%       \begin{definition}[2] …… \end{definition} → 「定义 2. ……」
\NewDocumentEnvironment{theorem}{o}
  {\par\medskip\noindent\textbf{定理\IfValueT{#1}{\ #1}.}\itshape\ignorespaces}{\par\medskip}
\NewDocumentEnvironment{proposition}{o}
  {\par\medskip\noindent\textbf{命题\IfValueT{#1}{\ #1}.}\itshape\ignorespaces}{\par\medskip}
\NewDocumentEnvironment{lemma}{o}
  {\par\medskip\noindent\textbf{引理\IfValueT{#1}{\ #1}.}\itshape\ignorespaces}{\par\medskip}
\NewDocumentEnvironment{corollary}{o}
  {\par\medskip\noindent\textbf{推论\IfValueT{#1}{\ #1}.}\itshape\ignorespaces}{\par\medskip}
\NewDocumentEnvironment{conjecture}{o}
  {\par\medskip\noindent\textbf{猜想\IfValueT{#1}{\ #1}.}\itshape\ignorespaces}{\par\medskip}
\NewDocumentEnvironment{claim}{o}
  {\par\medskip\noindent\textbf{断言\IfValueT{#1}{\ #1}.}\itshape\ignorespaces}{\par\medskip}
\NewDocumentEnvironment{definition}{o}
  {\par\medskip\noindent\textbf{定义\IfValueT{#1}{\ #1}.}\ignorespaces}{\par\medskip}
\NewDocumentEnvironment{remark}{o}
  {\par\medskip\noindent\textbf{注\IfValueT{#1}{\ #1}.}\ignorespaces}{\par\medskip}
\NewDocumentEnvironment{example}{o}
  {\par\medskip\noindent\textbf{例\IfValueT{#1}{\ #1}.}\ignorespaces}{\par\medskip}
\NewDocumentEnvironment{exercise}{o}
  {\par\medskip\noindent\textbf{习题\IfValueT{#1}{\ #1}.}\itshape\ignorespaces}{\par\medskip}
% 证明：直接用 amsthm 的 proof 环境，仅把提示词改为中文「证明」
\renewcommand{\proofname}{证明}

% 原文行内粗体小标题（"**Tame punctures.**"）→ 独立成段粗体
\newcommand{\para}[1]{\par\medskip\noindent\textbf{#1}\quad\ignorespaces}

% 图标题 / 表标题（手工编号）
\newcommand{\figtitle}[2]{\addcontentsline{lof}{figure}{\protect\numberline{#1}#2}\caption*{\small\textbf{图 #1：}#2}}
\newcommand{\tabtitle}[2]{\caption*{\small\textbf{表 #1：}#2}}

% 参考文献条目（生成脚本按 \refitem{编号}{内容} 输出）
\newcommand{\refitem}[2]{\par\hangindent=1.5em\hangafter=1\noindent [#1]~#2}

% 方框公式（原文若干编号公式是带边框的"盒子"，如式 (1.2)）
% 用法：\begin{equation}\eqbox{内容} \longleftrightarrow \eqbox{内容} \tag{1.2}\end{equation}
\newcommand{\eqbox}[1]{\fbox{\begin{minipage}{0.34\textwidth}\centering\small #1\end{minipage}}}

% ---------- 常用记号（只用这里定义过的宏，不要自造 \newcommand） ----------
\newcommand{\R}{\mathbb{R}}
\newcommand{\C}{\mathbb{C}}
\newcommand{\Z}{\mathbb{Z}}
\newcommand{\N}{\mathbb{N}}
\newcommand{\Q}{\mathbb{Q}}
\newcommand{\cA}{\mathcal{A}}
\newcommand{\cB}{\mathcal{B}}
\newcommand{\cC}{\mathcal{C}}
\newcommand{\cD}{\mathcal{D}}
\newcommand{\cE}{\mathcal{E}}
\newcommand{\cF}{\mathcal{F}}
\newcommand{\cG}{\mathcal{G}}
\newcommand{\cH}{\mathcal{H}}
\newcommand{\cI}{\mathcal{I}}
\newcommand{\cJ}{\mathcal{J}}
\newcommand{\cK}{\mathcal{K}}
\newcommand{\cL}{\mathcal{L}}
\newcommand{\cM}{\mathcal{M}}
\newcommand{\cN}{\mathcal{N}}
\newcommand{\cO}{\mathcal{O}}
\newcommand{\cP}{\mathcal{P}}
\newcommand{\cQ}{\mathcal{Q}}
\newcommand{\cR}{\mathcal{R}}
\newcommand{\cS}{\mathcal{S}}
\newcommand{\cT}{\mathcal{T}}
\newcommand{\cU}{\mathcal{U}}
\newcommand{\cV}{\mathcal{V}}
\newcommand{\cW}{\mathcal{W}}
\newcommand{\cX}{\mathcal{X}}
\newcommand{\cY}{\mathcal{Y}}
\newcommand{\cZ}{\mathcal{Z}}
\newcommand{\cHm}{\mathcal{H}}
\newcommand{\fg}{\mathfrak{g}}
\newcommand{\fh}{\mathfrak{h}}
\newcommand{\ft}{\mathfrak{t}}
\newcommand{\fu}{\mathfrak{u}}
\newcommand{\fgl}{\mathfrak{gl}}
\newcommand{\fsl}{\mathfrak{sl}}
\newcommand{\fsu}{\mathfrak{su}}
\newcommand{\fso}{\mathfrak{so}}
\newcommand{\fsp}{\mathfrak{sp}}
\newcommand{\dd}{\mathrm{d}}
\newcommand{\ee}{\mathrm{e}}
\newcommand{\ii}{\mathrm{i}}
\newcommand{\Tr}{\operatorname{Tr}}
\newcommand{\tr}{\operatorname{tr}}
\newcommand{\diag}{\operatorname{diag}}
\newcommand{\sgn}{\operatorname{sgn}}
\newcommand{\id}{\mathrm{id}}
\newcommand{\rk}{\operatorname{rk}}
\newcommand{\vol}{\operatorname{vol}}
\newcommand{\Pf}{\operatorname{Pf}}
\newcommand{\Res}{\operatorname{Res}}
\newcommand{\Mor}{\operatorname{Mor}}
\newcommand{\Ob}{\operatorname{Ob}}
\newcommand{\Hom}{\operatorname{Hom}}
\newcommand{\End}{\operatorname{End}}
\newcommand{\Aut}{\operatorname{Aut}}
\newcommand{\Bun}{\operatorname{Bun}}
\newcommand{\Rflat}{\mathcal{M}_{flat}}
\renewcommand{\Re}{\operatorname{Re}}
\renewcommand{\Im}{\operatorname{Im}}
\newcommand{\half}{\tfrac{1}{2}}
\newcommand{\ie}{\textit{i.e.}}
\newcommand{\eg}{\textit{e.g.}}
\newcommand{\cf}{\textit{cf.}}
\newcommand{\dbar}{\bar{\partial}}
\newcommand{\abs}[1]{\left|#1\right|}
\newcommand{\norm}[1]{\left\|#1\right\|}
\newcommand{\ip}[2]{\langle #1,#2\rangle}
\newcommand{\LG}{{}^{L}G}
\newcommand{\LT}{{}^{L}T}
\newcommand{\Lg}{{}^{L}\mathfrak{g}}

% hyperref 书签里遇到数学（\section 标题含 $...$）时不报错
\pdfstringdefDisableCommands{%
  \def\LG{LG}\def\C{C}\def\R{R}\def\Z{Z}\def\ii{i}\def\dd{d}%
}
